\section{Implementation}

MeshFlow is implemented in Go (2,873 lines of core logic) and ships as a single statically linked binary. Table~\ref{tab:code-size} breaks down the implementation by component.

\begin{table}[h]
\centering
\caption{Lines of code by component.}
\label{tab:code-size}
\begin{tabular}{lr}
\toprule
\textbf{Component} & \textbf{LOC} \\
\midrule
DAG Engine (orchestration, claiming, retry) & 537 \\
Web Dashboard (HTML/CSS/JS SPA) & 485 \\
CLI Entrypoint & 339 \\
Node Composer (relay, heartbeat, lifecycle) & 331 \\
HTTP API Server + SSE & 213 \\
SWIM Gossip Wrapper & 189 \\
BoltDB Persistence & 174 \\
HTTP Executor + Retry/Backoff & 167 \\
Embedded NATS + JetStream & 157 \\
Workflow/Task/Run Types & 142 \\
DAG Parser + Cycle Detection & 131 \\
Configuration & 123 \\
Event Types & 61 \\
\midrule
\textbf{Total Core Logic} & \textbf{2,873} \\
\bottomrule
\end{tabular}
\end{table}

The binary is built with \texttt{CGO\_ENABLED=0}, producing a fully static 26 MB artifact (20 MB in a multi-stage Docker Alpine image) with zero runtime dependencies. No external database, message queue, or coordination service is required. The web dashboard---a 485-line single-page application with SVG DAG rendering, mesh topology visualization, and live SSE event streaming---is embedded via Go's \texttt{//go:embed} directive and served directly from the binary.
